\chapter{Permission before value}\label{ch:permission}

\section{A positive number cannot create a physical possibility}
An EBU calculation can evaluate a hypothetical endpoint. It cannot make that endpoint physically attainable. A source with three units cannot supply an isolated four-unit lossless transfer. A route that does not exist cannot become available because its imagined value is high. A tank's ownership, access rules, and equipment limitations are also not determined by a potential difference.

Permission therefore comes before treating a quote as an executable offer. The phrase does not mean that every physical permission rule has already been settled. It means the record must identify a declared admissible action before a value can promise anything about its accepted execution.

The first model represents internal relocation of a conserved scalar. It does not import a regenerative extraction rule merely because both rules use source stocks. Moving resource between represented stores and withdrawing a resource from a regenerating ecological stock are different physical questions. Their permissions must be justified in their own domains.

\section{Three questions with three answers}
The physical question asks whether the action can occur under the declared world and event rules. The valuation question asks how it changes the fixed potential, with any justified process burden. The account question asks whether a proposed owner or actor can afford the signed change under a separately declared capacity policy.

A legal or ethical permission question can add another layer in an application. This textbook does not collapse it into physical feasibility. A mechanically feasible transfer can still be unauthorized by the people responsible for the resource. Conversely, institutional approval cannot supply missing physical stock.

\begin{center}
\begin{tabular}{P{0.25\linewidth}P{0.62\linewidth}}
\toprule
Layer & Example of its question\\
\midrule
Physical feasibility & Does the source contain the required stock, and is the declared route available?\\
Finite valuation & Does the accepted transformation lower or raise the declared potential?\\
Capacity policy & Does the signed account update satisfy the proposed balance rule?\\
Institutional authority & Who may authorize the action, and what rights or duties apply?\\
\bottomrule
\end{tabular}
\end{center}
The layers can interact through a specified protocol. Their interaction does not make them the same object. Keeping their answers separate is especially important when an action is refused: one should be able to say which rule failed.

\section{The simplest isolated transfer domain}
For an isolated lossless edge \(i\to j\) with no additional storage bound, no intervening flow, and nonnegative stock, a finite quantity must satisfy
\begin{equation}
 0\leq q\leq z_i.
\end{equation}
The lower bound respects the declared direction. The upper bound ensures the source successor \(z_i-q\) remains nonnegative. The destination gains a nonnegative quantity, and total stock remains unchanged.

This interval is the physical domain of that simple event. In the main fixture it is \([0,18]\). Every point in it is physically possible under these particular assumptions, even though its EBU may be negative. The \(q=17\) point is feasible and has field value \(-4.25\). The \(q=19\) point is outside the domain and cannot be accepted as this transfer.

Real infrastructure may impose a smaller maximum quantity, a finite receiving capacity, a time-dependent access restriction, or a quality condition. Those limits must be stated as physical inputs. The Gaussian scale is not a substitute for them. A convenient potential value cannot waive them.

\section{One source, several requests}
Suppose Ridge holds eight units and two actors each propose to send five units to different destinations. Each request alone is no larger than eight. Together they require ten. If the event supplies both from the same frozen source without incoming stock available during execution, the group is infeasible.

The group must be tested as a group. For that declared event semantics, a source-availability condition is
\begin{equation}
 \sum_{a:\,\mathrm{src}(a)=i}q_a\leq z_i.
\end{equation}
This condition is not an allocation rule. It says which sets of quantities pass the source check. Choosing which actor receives a reduced quantity, whether all requests are proportionally scaled, or whether a different group is selected requires a further declared mechanism.

There is no scientific need to pretend those choices are determined by conservation. Conservation constrains the total; it does not decide fairness, priority, or actor sampling. A later protocol can choose among alternatives before outcomes are seen and test the resulting process honestly.

\section{Incoming stock and event semantics}
Consider a three-cell chain. The middle cell begins empty. One action sends four units into it, while another sends four out. The net endpoint leaves the middle cell empty, with no negative stock anywhere. Does that prove the group is physically feasible?

Only if the declared physical event permits this synchronized passage. A continuous pipe can transmit incoming flow while releasing outgoing flow, provided its complete physical constraints are satisfied. A warehouse dispatch protocol may require goods to be received and verified before they can be sent onward. In that second model, an empty middle store cannot fund the outgoing action at the same frozen event.

The distinction cannot be resolved by adding or subtracting the same columns. The endpoint calculation is common to both stories. A protocol must additionally specify the realized path or availability rule. Until then, the right answer is that endpoint feasibility has been checked while event feasibility remains unspecified.

\begin{cautionbox}{A necessary check is not a complete permission}
Nonnegative endpoints and conserved totals are necessary in the first domain. They may be insufficient when a group depends on timing, shared equipment, transient availability, or other declared physical constraints. A complete acceptance rule must name those constraints instead of inferring them from the final sum.
\end{cautionbox}

\section{Joint feasibility should not be smuggled into valuation}
Suppose a group has an attractive joint field reduction but violates a source limit. One cannot rescue it by dividing its value among actors differently. The physical quantities are infeasible before any attribution is considered.

Conversely, a feasible group may have a joint value that differs from the sum of independent quotes. That is a valuation issue after the physical increment is established. The remedy is to calculate the joint field value from its common baseline and endpoint, then apply a declared attribution convention if one is needed. It is not to modify stocks until a preferred sum appears.

This separation makes debugging and scientific review more precise. A rejected group can fail because it double-used stock, because its receipt split was undefined, or because an owner could not afford a negative attribution. These are different propositions, each with different evidence and remedies.

\section{Candidate menus are restrictions, not all of physics}
The scalar resource is continuously divisible in the ideal model. A finite candidate menu might nevertheless offer only quantities one, two, and four on an edge. That restriction can make a future experiment finite and reproducible. It does not prove that intermediate quantities are physically impossible.

If the best recovery quantity is three and the menu offers only one and four, an observed outcome can depend on the menu. A comparison must state that limitation. Changing the menu after seeing that EBU performs poorly would change the registered problem. So would including many duplicates of one action if selection is uniform over menu entries.

The menu must also distinguish a no-action option from an empty feasible set. A nonempty menu containing rest can select rest by policy. An empty menu may trigger a separately declared no-action event. Both leave the physical state unchanged in the simplest case, but they have different explanations and can have different probabilities in a stochastic process.

\section{Permission is not a preference ranking}
An admissible set tells us which options are allowed. A selector chooses among them. The set may contain actions with high, low, zero, and negative EBU. Removing every negative action would add a sign-based filter. Choosing the largest positive action would add an optimizer. Neither is implicit in the definition of physical feasibility.

The prospective programme proposes an EBU-independent construction of candidates, joint physical checks, exact value and attribution, owner affordability, and then unweighted random selection among the surviving options. This order is a design direction with remaining protocol decisions. It must not be rewritten in explanatory prose as ``the field chooses the best action.''

Why random selection? It makes a particular research question possible: whether a declared signed capacity mechanism changes behaviour without directly rewarding the actor with an optimizer that always picks the largest field improvement. Whether that mechanism is useful remains open. The potential supplies values; the stochastic rule determines how those values affect available choices.

\section{The prospective capacity boundary}
The programme describes nonnegative persistent balances and signed updates assigned through a declared ownership map. For an affected owner \(i\), an affordability check would require
\begin{equation}
 B_i+\Delta B_i\geq0.
\end{equation}
Here \(B_i\) is proposed action capacity, not a physical stock and not a legal monetary entitlement. The update \(\Delta B_i\) depends on the separately specified group attribution and owner mapping. Its meaning does not follow from the symbol \(B\) used in another historical context.

For an illustrative arithmetic check, an owner with balance five and proposed update minus four remains at one. With update minus six, the proposed event fails that balance condition. Neither calculation changes the physical water in Ridge. If a transfer is physically impossible, a large balance cannot make it possible.

The fact that child attributions sum to a group field reduction does not determine who owns them or whether simultaneous positive and negative entries may net within one owner. Those are policy choices. The prospective programme identifies them as decisions for registration; this chapter does not adopt them merely by presenting the inequality.

\section{Why preservation of every local stock is the wrong question here}
An internal transfer deliberately reduces its source stock and increases its destination stock. Requiring every source to remain unchanged would eliminate every nonzero transfer. That may be appropriate for a particular protected resource or a distinct extraction policy, but it is not a consequence of global mass conservation.

Our Gaussian reference can value source reduction positively when the source is above its reference. It can also value it negatively when it is below. The value is an account of deviation, not an extraction permission. A hard rule protecting a source for some domain reason must be declared separately and assessed for that domain.

This distinction avoids reusing a historically valid theorem outside its assumptions. A regenerative source certificate can depend on growth, reserves, and a no-export successor. The closed internal-transfer world contains none of that regenerative law. Importing the certificate's arithmetic while omitting its eligibility assumptions would not inherit its authority.

\section{What a permission record should expose}
A reader should be able to reconstruct the proposed event from the permission record. It names the state epoch, physical action identifiers, finite quantities, shared resources, relevant physical limits, and the acceptance result. If accepted quantities differ from requested quantities, the mechanism producing that change must be stated.

The record also identifies what was not checked. If it verifies only endpoints, it must not claim a time-resolved physical path. If it uses a menu restriction, it must not claim exhaustive continuous feasibility. If a resource-quality issue lies outside the model, it should remain visible as a limitation.

For a real application, measurement provenance and institutional authority would accompany these data. In our textbook arithmetic, the inputs are declared illustrative values. The structure is still worth learning because it makes the difference between a physical statement and an inferred permission explicit.

\section{Guided problem: four candidate actions}
At the main baseline, consider the candidate quantities \(0,4,17,19\) on Ridge-to-Vale. Under the simple isolated transfer rule, zero, four, and seventeen are feasible; nineteen is not. Their field values are respectively zero, twelve, and minus 4.25. The value of nineteen can be evaluated as a formal polynomial, but it is not an admissible schedule point for this event.

Now add, for this exercise only, an independently declared route quantity limit of six. The feasible listed candidates become zero and four. The cap comes from the hypothetical route declaration, not from the potential. Do not change the field formula when applying it.

Finally suppose a separately declared capacity policy requires the actor to fund negative signed changes and the actor's capacity is three. Under the original eighteen-unit physical limit, the seventeen-unit action would be physically feasible but fail this particular affordability check. With capacity five it would pass the illustrative account condition. Its physical quantities and its EBU remain the same in both cases.

\section{Guided problem: two actors share eight units}
Source stock is eight. Actor A requests five for destination A, and actor B requests five for destination B. Three possible groups are \(\{A\}\), \(\{B\}\), and \(\{A,B\}\), with the requested quantities fixed. Under a no-incoming source-availability rule, the single-action groups pass and the two-action group fails.

One possible future allocation mechanism could reduce each request to four. Another could choose one request unchanged. A third could reject the contested group and select from the passing groups. These mechanisms produce different physical increments and potentially different field values. No choice is adopted by this example.

A correct comparison would define the mechanism before computing and interpreting outcomes. If a subset resolver changes quantities, its subset value must be calculated from the quantities that resolver actually declares. It cannot be substituted into a fixed-increment interaction formula without acknowledging the change of object.

\section{Failure is useful information}
Refusal need not mean the theory has broken. It can mean a physical constraint is doing exactly what it says. The scientific question is whether the permitted process can achieve useful outcomes under its declared conditions, and how its failures compare with a suitable alternative.

A permanently empty menu can also expose a badly posed study. If every action is prohibited by an imported policy unrelated to the new domain, no accounting mechanism is being meaningfully tested. Similarly, a world with no disturbances and every stock already at reference can make a zero-genesis affordability process inert. These are design facts to identify before attributing success or failure to EBU generally.

The remedy is to state the research question and the physical world consistently. It is not to choose permissions after inspecting outcomes until a favoured method acts. Prospective design can be revised openly before registration; reported evidence must retain the protocol under which it was produced.

\section{Practice and reflection}
A source begins with six units and two simultaneous outgoing quantities are four and three. Under the frozen-source rule used in the exercise, aggregate withdrawal seven fails even though each request passes alone. If an independent incoming transfer supplies two during the same event, explain why its presence still does not settle feasibility until timing or flow-through semantics are declared.

For an isolated main-fixture transfer of twelve units, the endpoint is \((6,14)\) and field reduction is twelve. It is physically feasible despite passing the reference. Explain why the fact that eight units would maximize the field reduction does not make twelve physically forbidden.

Finally list the information needed to diagnose a refusal. A complete answer distinguishes absent edge, insufficient stock, joint resource conflict, menu omission, undefined attribution, and failed affordability. Treating all these as ``the EBU says no'' would conceal the actual mechanism.
